\addcontentsline{toc}{subsubsection}{Solution 56. The Microscopic Error Trap}

\begin{dlsummary}[title={\bfseries Solution 56. The Microscopic Error Trap},breakable]

\begin{proof}
The complete family is
\[
\boxed{f(x)=ax\qquad(x\in\mathbb R),\quad a\in\mathbb R.}
\]

\paragraph{1. Extract exact doubling.}

Taking \(x=y=0\) in (1) gives \(f(0)=0\).

Taking \(y=x\), the upper bound vanishes, so
\[
\boxed{f(2x)=2f(x).} \tag{3}
\]
Consequently,
\[
f(x/2^n)=\frac{f(x)}{2^n}
\]
for every nonnegative integer \(n\).

\paragraph{2. Shrink the error until it disappears.}

Apply (1) at \(x/2^n,y/2^n\). Using (3),
\[
0\le
\frac{f(x)+f(y)-f(x+y)}{2^n}
\le\frac{(x-y)^2}{4^n}.
\]
Multiplying by \(2^n\),
\[
0\le f(x)+f(y)-f(x+y)
\le\frac{(x-y)^2}{2^n}.
\]
Letting \(n\to\infty\), we obtain
\[
\boxed{f(x+y)=f(x)+f(y).} \tag{4}
\]

Additivity alone is \textbf{not} enough to conclude that \(f\) is linear.

\paragraph{3. Find regularity in the composite.}

Define
\[
h(x)=f(f(x)).
\]
By (4), \(h\) is additive. Condition (2) says
\[
h(x^2)\ge f(x)^2\ge0.
\]
Every nonnegative real number is a square, so
\[
h(t)\ge0\qquad(t\ge0).
\]
Thus \(h\) is nondecreasing: for \(u<v\),
\[
h(v)-h(u)=h(v-u)\ge0.
\]

An additive, nondecreasing function is linear. Indeed, writing \(c=h(1)\ge0\), additivity gives \(h(q)=cq\) for rational \(q\), and rational approximation plus monotonicity gives
\[
\boxed{h(x)=cx\qquad(x\in\mathbb R).} \tag{5}
\]

\paragraph{4. Transfer the regularity back to \(f\).}

Combining (2) and (5),
\[
f(x)^2\le h(x^2)=cx^2.
\]
Therefore
\[
\boxed{|f(x)|\le\sqrt c\,|x|.} \tag{6}
\]
By additivity,
\[
|f(u)-f(v)|=|f(u-v)|\le\sqrt c\,|u-v|.
\]
Hence \(f\) is continuous.

Put \(a=f(1)\). Additivity gives \(f(q)=aq\) for every rational \(q\); continuity extends this to every real input:
\[
f(x)=ax.
\]

\paragraph{5. Verify every slope.}

For any real \(a\), the function \(f(x)=ax\) has zero additive error, and
\[
f(f(x^2))=a^2x^2=f(x)^2.
\]
Thus every real slope works, including \(a=0\) and negative slopes.
\end{proof}

\paragraph{Takeaways.}
\emph{No takeaways were provided in the source draft.}

\end{dlsummary}
